% !TEX root = ../main.tex
% ======================================================================
% Section 6 – Gestion, responsabilité et propriété intellectuelle
% Responsable : Alex D'Andrea
% Grille : « PI, gestion des données & responsabilité » – 14 pts (section prioritaire).
% Longueur indicative : ~3,5 à 4 pages.
% ======================================================================
\subsection{Propriété intellectuelle et responsabilité}
\label{sec:pi-responsabilite}
Cette section traite de la propriété intellectuelle et des responsabilités de chaque acteur.
Dans un modèle fédéré, la contribution est partagée entre plusieurs établissements, ce qui rend la question de la \ac{PI} plus complexe.

\subsubsection{Propriété du modèle global}
\label{sec:propriete-modele}
Le code source, le protocole de traitement des données et le nom \fedesante{} appartiennent à l'équipe qui les a conçus.
Le modèle, en soi, pose une autre question.
Sa valeur réside dans ses poids : sans eux, l'architecture est inutile.
Or, ces poids ont été ajustés à partir des données des 12 établissements participants.
\arediger{confirmer l'algo d'agrégation dans section metho}
Avec \ac{FedAvg} pondéré, les grands sites pèsent davantage dans l'agrégation, mais chaque contribution reste indispensable au modèle final.
Les auteurs proposent donc une copropriété du modèle global, gérée par un consortium.
Chaque établissement y disposera d'une voix, quelle que soit sa taille, afin d'assurer l'équité dans la gestion du modèle.
Cette copropriété ne suffit toutefois pas pour les trois centres de santé des Premières Nations, dont les données relèvent d'une gouvernance propre, abordée à la \cref{sec:pcap-parametres}.


\subsubsection{Principes PCAP appliqués aux paramètres du modèle}
\label{sec:pcap-parametres}
% TODO : comment appliquer les principes \pcap{} (propriété, contrôle, accès, possession) à des paramètres de modèle plutôt qu'à des données brutes ? (citer \cite{pcap})
\arediger{PCAP et paramètres}

\subsubsection{Retrait d'un site et désapprentissage}
\label{sec:retrait-desapprentissage}
% TODO : un site peut-il se retirer et exiger le « désapprentissage » (machine unlearning) de sa contribution ? Faisabilité technique et engagements contractuels.
\arediger{retrait et désapprentissage}

\subsubsection{Protocole d'agrégation : brevet ou code ouvert ?}
\label{sec:brevet-code-ouvert}
% TODO : le protocole d'agrégation doit-il être breveté ou publié en code ouvert ?
\arediger{brevet ou code ouvert}

\subsection{Documentation}
\label{sec:documentation}

\subsubsection{Fiche de modèle (\emph{model card})}
\label{sec:model-card}
% TODO : contenu de la fiche de modèle.
\arediger{fiche de modèle}

\subsubsection{Fiche de données par site (\emph{datasheet})}
\label{sec:datasheet}
% TODO : contenu de la fiche de données par site.
\arediger{fiche de données}

\subsubsection{Registre des participants}
\label{sec:registre}
% TODO : contenu et tenue du registre des participants.
\arediger{registre}

\subsubsection{Règles d'accès au modèle pour les chercheurs externes}
\label{sec:acces-externe}
% TODO : règles d'accès au modèle pour les chercheurs externes (académiques ou privés).
\arediger{règles d'accès}
